% =============================================================================
%  Lower Volume, Articles 3–42
%  Source: translation/pages_63-66.md, pages_67-70.md, pages_71-74.md,
%          pages_75-78.md
% =============================================================================

% ---------------------------------------------------------------------------
%  Article 3
% ---------------------------------------------------------------------------
\article{3}{When One Missteps, Everything Is Reversed}

\begin{sokyubox}
{\small\itshape\color{sokyubar} \jp{第三章\quad 一ト足踏みはづす時は皆裏表になる事}}

\medskip
When a trade has gone well for several months and one's forecasts
have proven correct, one must never ride the winning streak. One's
sole aim should be to close out safely. Because one's forecasts
have been right for months, one grows complacent --- both in buying
and in selling --- and at every small fluctuation one is tempted to
stake another view. Guard against this. In such a trade, the moment
one missteps, everything is turned inside out. Take profits
carefully, again and again; never treat the market carelessly. Rise
and fall in this trade are matters of natural order, and apart from
the Three-Position Tradition (\jt{三位の傳}{sanmi no den}) there is no
way to secure profit.
\end{sokyubox}

\commentary
When one has identified the direction months in advance, entered a
trade, and seen it bear fruit just as expected, one must not, drunk
on success, press for still more. The essential thing is simply to
take profit without incident (\jt{利喰}{rigui}). After months of
correct forecasting, one comes to regard the market as easy to read,
and at every small price movement one is eager to stake large
speculative positions. This must be strongly guarded against. To
form a market view is not so simple a matter. While one is winning,
all is well; but should one slip even once, everything inverts ---
the profits of recent months turn into losses of equal magnitude.
Therefore one must, again and again, study the footprint chart
closely, weigh the price level of the ceiling and floor, and trade
accordingly. Apart from the contemplative discipline of Zen
(\jt{禪心}{zenshin}), profit is hard to come by.

% ---------------------------------------------------------------------------
%  Article 4
% ---------------------------------------------------------------------------
\article{4}{The Secret Tradition of Turning One's Sentiment in Trading}

\begin{sokyubox}
{\small\itshape\color{sokyubar} \jp{第四章\quad 賣買買氣を轉じ秘傳の事}}

\medskip
When rice appears weak and selling pressure mounts urgently, wait
three days, reverse your sentiment, and join the buy side. This is
invariably profitable. When buying pressure mounts and everyone is
certain prices must rise, reverse your sentiment once more and sell.
This is the innermost secret (\jt{極意}{gokui}) of the rice trade ---
never forget it. When you yourself feel bullish, know that everyone
else feels the same; when you yourself feel bearish, everyone else
has likewise tilted bearish. What has risen to the limit falls; what
has fallen to the limit rises --- this is the natural principle of
yin and yang, and requires no further deliberation. Entrust yourself
entirely to the Three-Position Tradition.
\end{sokyubox}

\commentary
Near the floor, all traders lean bearish, and one's own impulse is
to sell along with them. But rather than selling at once, step back,
reflect for three days, reverse one's sentiment, and turn to the
buy side. The market will indeed rise. Conversely, near the ceiling,
people proclaim that there is no telling how much higher prices may
go; on the surface, everything that appears is bullish; one is
swept along by the crowd's enthusiasm and wants to buy. Yet here
too, one should make an abrupt about-face and turn to selling. The
market, you see, always runs ahead of the crowd's thinking. When
every last person has turned bullish and buys recklessly, the
market's hidden undercurrent has already shifted toward decline.
Rising to the limit and then falling, falling to the limit and then
rising --- this is the market's normal state and its fundamental
principle. Rather than being made a fool of by overheated sentiment,
one should maneuver according to the Three-Position Tradition taught
in this book.

% ---------------------------------------------------------------------------
%  Article 5
% ---------------------------------------------------------------------------
\article{5}{Guidance for a Year-Block Year When It Turns toward the West}

\begin{sokyubox}
{\small\itshape\color{sokyubar} \jp{第五章\quad 三年塞西の方へ廻る年心得の事}}

\medskip
In a year when the three-year directional taboo (\jt{年塞}{toshi-fusagari})
turns toward the west --- that is, the Boar, Rat, and Ox years
(\jp{亥子丑}) --- rice rises greatly. This principle admits
no exception. However, if government storehouse rice
(\jt{御蔵米}{okura-mai}) lies unsold in surplus, or disputes arise over
warehouse certificates (\jt{札方}{fudakata}), or bumper harvests have
continued for several years running and rice is plentiful at rock-bottom
prices, then hold back and do not form deep speculative positions.
If none of these adverse conditions apply, and the harvest across
Japan and in the local region is normal, and the price stands cheap
relative to the year's character, buy aggressively. In years when
the three-year taboo is in effect, natural disasters strike Kyushu,
the Kamigata region, and the local area alike --- do not doubt this.
After the ceiling price emerges during winter, there is a slight
pullback; thereafter the first-stem month (\textit{kinoe-tsuki}) falls in
the eleventh, first, or third month, and at that time prices rise
fifty or sixty tawara. Do not doubt it. In a year when the
first-stem month falls in the seventh month, even if Japan and the
local region both enjoy good harvests, do not hesitate --- from the
sixth month onward, buy without letting down your guard.
\end{sokyubox}

\commentary
This chapter is not necessarily to be taken on faith, but here the
commentator provides an explanation following the master's original text.

The three-year western taboo (\jp{西三年の塞}) listed in the calendar
apparatus at the end of this volume --- that is, the three years
Boar, Rat, and Ox --- are years in which the market rises greatly.
Based on past experience, this holds true in eight or nine cases out
of ten. Even so, if the deliverable rice --- that is, the warrant
rice (\jt{券米}{kenmai}) --- is plagued by empty-certificate issuance
and non-delivery, or if years of consecutive bumper harvests have
left surplus rice flooding the market, then even in a western-taboo
year the rise will be modest, and one should not form speculative
positions. If, however, no such adverse conditions exist, and the
harvest across Japan --- especially in the local region --- is good,
and the price stands cheap relative to the crop, one may speculate
fully on the buy side. Past precedent shows that in years when the
three-year taboo is in effect, the cycle brings natural disasters in
its train. After the winter ceiling emerges and prices pull back
slightly, the first-stem month (\textit{kinoe}) aligns with the eleventh,
first, or third month, and at each of these junctures prices rise
fifty or sixty tawara.

As for the great rise described above: in a year when the monthly
stem-branch rotation places the first-stem month in the seventh
month, one should turn to the buy side from the sixth month onward,
buying without letup. The expression ``government storehouse rice
lies unsold'' (\jp{御蔵米滯り}) in the proviso refers to stockpiled rice
that refuses to diminish.

% ---------------------------------------------------------------------------
%  Article 6
% ---------------------------------------------------------------------------
\article{6}{On Days When Nakagami Is in the Heavens}

\begin{sokyubox}
{\small\itshape\color{sokyubar} \jp{第六章\quad 天一天上日の事}}

\medskip
The days called \jt{天一天上}{ten'ichi tenj\=o},
\jt{八専入}{hassen-iri}, \jt{丑日}{ushi no hi},
and \jt{不成就日}{f\=uj\=oju-bi}:

When rice has been rising steadily and nears a hundred-tawara gain,
and the high price falls on one of these days, that price is the
annual ceiling (\jp{年中行附天井直段}). Thereafter, no higher price will
appear. In a three-year taboo year or a year when the first-stem
month falls in the seventh month, the ceiling will without fail
come on one of these days. In ordinary years, it is said that the
ceiling may also fall on these days, but in practice this is rare
--- in ordinary years, the ceiling comes on an unremarkable day.
When rice has been rising steadily toward a hundred-tawara gain and,
on the twenty-ninth and thirtieth of the month, sentiment is
buoyant by a \textit{bu} or two, one should understand that the ceiling price
will appear the following month.
\end{sokyubox}

\commentary
This chapter records observations drawn from the three-year
western-taboo and seventh-month first-stem years. In such years,
when the market has risen to its limit, the ceiling will absolutely
come on one of the days listed above. In other years, however,
there is no fixed rule --- the ceiling falls on an ordinary day.
The way to read the ceiling is this: when the market has been
climbing steadily for several months and, around the end of the
month, the internal sentiment (\jt{内気配}{uchi-kehai}) is
buoyant, one should understand that the ceiling price will emerge the
following month.

% ---------------------------------------------------------------------------
%  Article 7
% ---------------------------------------------------------------------------
\article{7}{Guidance for When One Thinks to Profit Quickly}

\begin{sokyubox}
{\small\itshape\color{sokyubar} \jp{第七章\quad 急に儲く可しと思ふ時心得の事}}

\medskip
When one hastens to profit quickly, one is bewildered by the daily
fluctuations and ends up chasing the market from trade to trade,
suffering losses each time. For new-crop rice, carefully consider
the weather (\jt{順氣}{junki}), the crop conditions in the Kamigata
and other regions, the ceiling and floor, and check these against
the Three-Position Tradition; then join the buy side and maintain a
bullish stance throughout. Spring and summer rice follow the same
order. The initial entry (\jt{踏出し}{fumidashi}) is critical: wait for
months on end and identify the floor --- this is the first
priority. When a trade is going well, profit-taking (\jt{見切}{mikiri})
is critical. After closing a trade, resting for forty or
fifty days is essential.
\end{sokyubox}

\commentary
When one rushes to profit, one overestimates the price range, chases
the market, and buys at the very moment prices are about to turn
downward from the ceiling, or sells short at the very moment prices
are about to turn upward from the floor. Crashing into one wall and
then the other, reversing position (\jt{ドテン}{doten}) each time, one
suffers losses on every trade. The market should be traded on the
grand view --- selling at the ceiling and buying at the floor.
There is no surefire method for capturing the middle swings.
Consider the crop conditions across all regions, the state of the
weather, market sentiment, the volume of rice in storage and rice
coming to market; fix one's attention on the ceiling and the floor.
When one has identified the ceiling, sell single-mindedly; when one
has identified the floor, buy single-mindedly --- trade with
decisive conviction. The initial entry is of the utmost importance,
so wait for any number of months, watching for the ceiling or floor
to appear before committing. Then take profit at the seven- or
eight-tenths mark, close the trade, rest quietly for forty or fifty
days, observe the market's course at leisure, and enter the next
trade.

% ---------------------------------------------------------------------------
%  Article 8
% ---------------------------------------------------------------------------
\article{8}{The Root of Rise and Fall Is the Quality of the Harvest}

\begin{sokyubox}
{\small\itshape\color{sokyubar} \jp{第八章\quad 高下の本は作の善惡の事}}

\medskip
The quality of the harvest is the root of all rise and fall. Each
year, the first priority is to consider the volume of rice ---
large or small --- from Kyushu, the Kamigata region, the local area
and its neighboring provinces, as well as old-crop rice. The
Three-Position Tradition is not a technique for knowing rise and fall
in itself; it is a technique for recognizing, when the time of the
Three Positions arrives, what kind of development has emerged and
whether the moment is ripe for a rise. This deserves deep
reflection.
\end{sokyubox}

\commentary
Since the quality of the harvest --- that is, the size of the
crop --- is the fundamental cause of market rise and fall, one must
assess the surplus rice carried over from the previous year together
with the current year's crop. The meaning is clear enough, but a
point of interpretive difficulty arises with the term ``Three
Positions'' (\jt{三位}{sanmi}). This term, which has appeared repeatedly
as a technical term in earlier chapters, admits of three
interpretations:

\textit{First interpretation}: ``Three Positions'' refers to the position of
the ceiling, the position of the floor, and the position of the
middle price.

\textit{Second interpretation}: ``Three Positions'' refers to the three-year
western-taboo and first-stem month system.

\textit{Third interpretation}: this is the view of Master Kuzuoka
(\jp{葛岡翁}), offered in his \textit{Sanmi Mond\=o}
(\jp{三位問答}). He held that ``Three Positions'' is an abbreviation of
\textit{sanmai} (\jt{三昧}{sanmai}, ``meditative absorption''), so that
``the Three-Position Tradition'' means: with a Zen mind and with the
aid of the price-level charts, to observe the market's position and
touch the heartstring of its rise and fall. Master Kuzuoka, however,
acknowledged the validity of the first two interpretations as well,
leaving the question open. To restate the chapter's meaning briefly:
one should calculate and weigh the fundamental causes of rise and
fall described in the preceding chapters, reflect deeply, wait for
months on end, and only when it accords with one's own conviction
should one commit to a trade (\jt{商内}{akinai}). This is what the
chapter teaches.

% ---------------------------------------------------------------------------
%  Article 9
% ---------------------------------------------------------------------------
\article{9}{Maneuvering by the Osaka Market}

\begin{sokyubox}
{\small\itshape\color{sokyubar} \jp{第九章\quad 大阪相塲を以て掛引の事}}

\medskip
In the Osaka market (\jt{大阪相塲}{\=Osaka s\=oba}), autumn prices tend to
rise on the strength of the various provinces' harvests. From spring
and summer through the seventh and eighth months, when old rice is
running thin and the wheat harvest has shown no shortfall, prices
decline year after year from the fourth month through the eighth.

When the local market enters its summer-rice season, ships arrive
faster than the Osaka market's pace, domain storehouse rice
(\jt{御蔵米}{okura-mai}) thins out, and rice runs scarce around the
Inland Sea coast. When ships compete to buy, unexpectedly high
prices can emerge. At such times, the volume of storehouse releases,
the flow of domestic rice, and the number of ships are the primary
considerations. One should note, year by year, around what time
Osaka reaches its ceiling (\jt{天井}{tenj\=o}) and around what time it
reaches its floor price (\jt{底直段}{soko nedan}), and maneuver
(\jt{掛引}{kakehiki}) accordingly. In Kaga rice, forty-eight or
forty-nine \textit{monme}, or fifty-one or fifty-two \textit{monme}, marks the floor.
\end{sokyubox}

\commentary
This chapter describes how to use the Osaka market as a reference
for maneuvering in the local market.

The normal pattern for the Osaka rice market through the year is as
follows: in autumn, prices tend to rise on the strength of each
province's crop conditions. From spring and summer through the
seventh and eighth months, when there is no concern about old rice
running short and no failure in the wheat harvest, prices decline
year after year from about the fourth month through the eighth.

In the local market's normal pattern, once summer rice comes in, the
pace is faster than Osaka's: buying ships arrive earlier than usual,
local stocks diminish, the provinces around the Inland Sea run short
of rice, and ships compete to buy. In such years, unexpectedly high
prices can appear. At those times, the most important tactical
considerations are the volume of local warehouse stocks, the volume
of farmers' held rice, the volume of export rice, and the number
of buying ships. One should keep a record of when Osaka typically
makes its ceiling and floor each year and apply that knowledge to
local maneuvering. As a rule of thumb, Kaga rice at forty-eight or
forty-nine \textit{monme} to about fifty \textit{monme} generally marks the floor.

% ---------------------------------------------------------------------------
%  Article 10
% ---------------------------------------------------------------------------
\article{10}{On Price Ranges}

\begin{sokyubox}
{\small\itshape\color{sokyubar} \jp{第十章\quad 十餘二三丁一二三匁廿匁より其餘の事}}

\medskip
When rice has risen more than ten \jt{俵}{tawara} above
the floor, one should sell. However, in a year of poor harvest, one
should wait and observe.
\end{sokyubox}

\commentary
Rice that has risen ten or more tawara above the floor should be
sold against the trend. However, in years of poor harvest, prices
can rise as much as a hundred and eighty tawara, so one must observe
carefully before selling.

\begin{sokyubox}
When rice has fallen two or three increments (\jt{丁}{ch\=o}), one should
buy. However, rice that has been pushed up to around twenty-five
tawara should not be bought in haste.
\end{sokyubox}

\commentary
Even in a sharp decline, once rice has passed one full stage of
major decline and dropped through a second and third breath of
selling, one may try buying tentatively. However, when the market
stands at around twenty-five tawara --- a high price level --- one must
not rush to buy.

\begin{sokyubox}
One or two \textit{monme} to three \textit{monme} of spread marks the ceiling; one
should sell. Twenty \textit{monme} or more of spread marks the floor; one
should buy. However, even if one notes the ceiling and floor prices,
when rice is in the midst of rising, calculations alone will not
hold.
\end{sokyubox}

\commentary
This chapter describes a method for estimating the local ceiling
and floor by using the spread (\jt{鞘}{saya}) --- that is, the carry
rice (\jt{乘米}{nosei-mai}) or price differential --- between the Osaka
market (the demand center) and the local market.

The references in the text to ``one or two \textit{monme}'' and ``twenty \textit{monme}''
refer to the spread of Sh\=onai rice against the Osaka market. The
meaning is: when the spread narrows to one, two, or three \textit{monme},
the market is very likely at the ceiling, and one should sell. When
the spread has widened to as much as twenty \textit{monme}, the market is
at the extreme floor, and one should buy. However, this is stated
in broad terms. In a rising market, the spread can narrow to nothing
or even invert, so one cannot maneuver on the spread calculation
alone. One must take a broad view of the overall trend.

% ---------------------------------------------------------------------------
%  Article 11
% ---------------------------------------------------------------------------
\article{11}{On the Twenty-Monme Floor}

\begin{sokyubox}
{\small\itshape\color{sokyubar} \jp{第十一章\quad 二十匁底の事}}

\medskip
The twenty-\textit{monme} spread marks the floor. In a normal year, old
rice declines around the fifth and sixth months. When two successive
years of good harvest occur, the spread can widen to as much as
twenty-five or twenty-six \textit{monme}. In a year of local crop failure,
the floor comes at a spread of sixteen or seventeen \textit{monme}, and the
ceiling arrives when the spread vanishes entirely. In a year of
good harvest, one should reckon that the floor is at a twenty-\textit{monme}
spread and the ceiling at eleven or twelve \textit{monme}. When the crop
outlook is uncertain, one should buy without hesitation. After the
ceiling price has appeared, one should gradually sell over five or
six months --- this yields full profit (\jt{十分利運}{j\=ubun riun}).
Until the ceiling price has appeared, one must absolutely not sell.
Note, however, that in the fifth and sixth months, depending on
the volume of rice, the situation may defy analysis. In years when
long positions are heavy, a collapse will certainly come in the
sixth month --- do not doubt it.
\end{sokyubox}

\commentary
This chapter elaborates in detail on the method introduced in
Article~10 for observing the ceiling and floor through the spread.

Around the floor, the spread against Osaka is typically about
twenty \textit{monme}. In a normal year, old rice declines around the fifth
and sixth months, and when two or three years of decent harvests
have continued, the spread can widen to twenty-five or twenty-six
\textit{monme} as prices pull back. In a year of local crop failure, the
floor comes at a spread of sixteen or seventeen \textit{monme}, and by the
time the ceiling price appears, the spread has vanished entirely. In
a year of good harvest, if the floor is at a twenty-\textit{monme} spread,
the ceiling comes at about eleven or twelve \textit{monme}. When the crop
outlook looks uncertain, one should buy at these spread levels
without hesitation. However, the fifth and sixth months are driven
so much by the volume of remaining stocks and long positions that
they are difficult to analyze with precision. In years when long
positions are heavy, a major break will certainly come around the
sixth month.

% ---------------------------------------------------------------------------
%  Article 12
% ---------------------------------------------------------------------------
\article{12}{On the Twelve-Monme Middle}

\begin{sokyubox}
{\small\itshape\color{sokyubar} \jp{第十二章\quad 十二匁中の事}}

\medskip
Twelve \textit{monme} of spread marks the middle price (\jt{中直段}{ch\=u nedan}).
When new rice first appears around the sixth or seventh month and
old rice is sitting at its fifth- or sixth-month ceiling price, new
rice will open at a correspondingly high price. At such times, one
should carefully assess the crop conditions and sell aggressively.
When the harvest is good, the spread can widen to more than twenty
\textit{monme}. However, with a stake of three hundred \textit{ry\=o}, quick-eyed
deep commitment is unnecessary --- exercise restraint.
\end{sokyubox}

\commentary
The ``middle'' (\jp{中}) means the middle price level. The reference
to ``three hundred \textit{ry\=o}'' and ``quick-eyed'' and so on means that one
should not commit large positions; rather, one should limit oneself
to small positions of two or three hundred \textit{ry\=o}, take profits
quickly, and cut losses promptly. The rest of the meaning is
sufficiently clear and requires no further explanation.

% ---------------------------------------------------------------------------
%  Article 13
% ---------------------------------------------------------------------------
\article{13}{On the One-Monme Ceiling}

\begin{sokyubox}
{\small\itshape\color{sokyubar} \jp{第十三章\quad 一匁天井の事}}

\medskip
One \textit{monme} of spread marks the ceiling. In a year of local crop
failure, the floor comes at a spread of sixteen or seventeen \textit{monme},
and one should reckon the ceiling at one \textit{monme} or with no spread at
all. From there, the spread widens: first to three or four \textit{monme},
then to eight or nine \textit{monme}, then down to twelve or thirteen \textit{monme}.
When this ceiling appears, one should study the spread and sell into
it. In a range-bound market (\jt{通相塲}{kayoi s\=oba}), prices will
assuredly decline for five or six months. In a falling market, one
should close out sell positions month by month. Yet when one has
sold forward positions two or three months out at the extreme
tightening of the spread with no carry at all, there is no need
to close out: \textit{``At two, close out; at three, fully ripe; at four,
it turns''} --- this is the secret of the Three-Position Tradition
(\jt{三位の傳}{sanmi no den}). Sell positions in an ordinary month
should be closed by the twenty-ninth. Sell positions in the fifth
and sixth months should be closed cleanly by the eighteenth,
nineteenth, or twenty-first. Bear this in mind.
\end{sokyubox}

\commentary
This chapter continues from the preceding one, describing the
behavior of the spread at the ceiling and the corresponding
maneuvering.

The meaning is as follows: at the ceiling, the spread narrows to
one \textit{monme} or to nothing. When the market has been rising for several
months and reaches this state of sentiment, it is assuredly the
ceiling. One should study the spread and sell in. The decline will
continue for two to three months, or as long as five to six months.
However, a falling market, as has been explained many times before,
drops in stages --- declining, then pausing, then declining again.
One should maneuver accordingly. If one's position is in the
near-month contract (\jt{期近物}{kijika-mono}), one should generally
close out at a reasonable level. The instruction on closing fifth-
and sixth-month sell positions has already been explained in the
preceding chapter and should be consulted there.

% ---------------------------------------------------------------------------
%  Article 14
% ---------------------------------------------------------------------------
\article{14}{Buying Yields Eight Parts Profit; Selling Yields Two}

\begin{sokyubox}
{\small\itshape\color{sokyubar} \jp{第十四章\quad 買八分の利賣二分の利と申事}}

\medskip
\textit{``Buying yields eight parts profit; selling yields two.''}

The seller collects interest and, when carrying profits on top,
appears to be in a most advantageous position. Yet if a crop
failure should strike the Kamigata region or the local area,
Osaka prices surge upward and buying back becomes impossible ---
losses result. The buyer, on the other hand, pays months of
interest and bears indeterminate costs of settlement, making
buying seem a most difficult proposition. Yet when fortune turns
favorable, the carrying costs become as dust before the wind and
cause no distress at all. Moreover, regardless of scale, autumn
rice never fails to rise --- one should not be a seller.
\end{sokyubox}

\commentary
This chapter declares that buying is the more advantageous side, and
expresses the proportion numerically: roughly eight parts to two in
favor of the buyer. It places great weight on the buy side.

That buying is the more profitable side is confirmed by the fact
that throughout history, the great majority of successful market
operators have been buyers. The author of this text, Master S\=oky\=u,
is of course an emphatic advocate of the buying side. It goes
without saying. But even Ushida Jiunsai (\jp{慈雲齋}) was
an early proponent of this view, declaring in the most extreme
terms: \textit{``Compare a lantern with a temple bell --- selling and
buying: buying profits more; selling loses more.''} He explained
the reason as follows: rice is by nature a commodity biased toward
strength, and therefore the advantage lies with the buyer. The
allies of bullish sentiment include wind, rain, drought, blight,
floods, and other natural disasters, as well as emergencies and
shipping disruptions. The allies of bearish sentiment, apart from
favorable weather and rice imports, have not a single supporting
factor. This chapter expresses the same principle.

The proviso explains the trading customs of the Sakata rice
exchange (\jt{酒田米会所}{Sakata kome kaisho}): when a seller does not
close out within the month in which the position was opened and
carries it over to the following month, the buyer must pay the
seller a monthly interest charge at the time of settlement. This
has already been explained in the Opening Volume. Thus, while the
seller receives interest from the buyer and, when the market moves
favorably, profits handsomely, should a crop failure occur, prices
surge sharply upward and buying back is delayed, resulting in
losses. By contrast, while the buyer must pay interest --- an
apparent disadvantage --- when the market turns favorable, the
profits dwarf the interest paid, which becomes as a single hair
among nine oxen (\jt{九牛が一毛}{ky\=ugy\=u ga ichim\=o}) and causes no
concern at all. Particularly in autumn, there is always talk of
insufficient sickle-cutting, overpraised crops, poor autumn
ripening, and one thing or another --- the market is invariably
biased upward, and the seller's position is disadvantageous. Beyond
this, there is a further disadvantage for the seller: at delivery,
the buyer need only prepare funds, but the seller must assemble
funds \textit{and} collect rice for delivery --- a double burden.

% ---------------------------------------------------------------------------
%  Article 15
% ---------------------------------------------------------------------------
\article{15}{Do Not Sell Rice in a Bountiful Year}

\begin{sokyubox}
{\small\itshape\color{sokyubar} \jp{第十五章\quad 豊年に米売り不可申事}}

\medskip
There is a saying: do not sell rice in a bountiful year. Because it
is a bountiful year, sentiment is weak and the market is priced
cheaply in anticipation. Yet a uniformly good harvest across all of
Japan is rare. Some province will produce a shortfall, and from the
ninth and tenth months, Osaka will begin to rise. Moreover, because
word of a bountiful year spreads everywhere, buyers will come even
through the winter.
\end{sokyubox}

\commentary
In a year that shapes up as bountiful, favorable conditions improve
steadily and the market prices in the good harvest by declining
gradually. By the time the bountiful year is confirmed, the market
has already reached its floor. However, a uniformly excellent
harvest across the whole of Japan is extremely rare; some province
inevitably produces a shortfall in crop conditions, and this serves
as the catalyst for prices to turn upward. This is the reason behind
the admonition: do not sell rice in a bountiful year.

% ---------------------------------------------------------------------------
%  Article 16
% ---------------------------------------------------------------------------
\article{16}{Do Not Buy Rice in a Famine Year}

\begin{sokyubox}
{\small\itshape\color{sokyubar} \jp{第十六章\quad 凶年に米買不可申事}}

\medskip
There is a saying: do not buy rice in a famine year. Because it is
a famine year, trade opens at high prices in anticipation. On top of
that, merchants from the local area and elsewhere pile on with heavy
speculative positions. In addition, townspeople and rural folk alike
buy even rice they do not need, what with tax-rice payments and
other causes, so that by the seventh, eighth, ninth, and tenth
months, prices are pushed to their limit beyond what the crop
shortfall alone would warrant. Furthermore, sentiment is stretched
to its extreme and held there through the summer. The innermost
truth is that prices will decline. Bear this in mind.
\end{sokyubox}

\commentary
This chapter is the reverse of the preceding one. In a famine year,
the market trades at an inflated premium. An inflated market will
inevitably decline in the end. Therefore, at the point where prices
have been pushed to their extreme, one should exercise great
caution and must not buy.

Now, in a famine year, local bullish sentiment (\jt{強氣}{tsuyoki}) goes
without saying, but speculators from other regions also crowd in with
buy positions, and the general public develops an urge to stockpile
even rice they have no immediate need for. The result is that the
market overshoots the actual crop shortfall by three or four months,
sometimes five or six, climbing steadily. Moreover, because people
witness the famine with their own eyes, they chase the rally and
sentiment grows taut. By the following spring, around the third or
fourth month --- sometimes stretching into the fifth month --- the
market holds at high prices in consolidation (\jt{保合}{mochai}). But
from around mid-fifth-month into the sixth month, it invariably
collapses.

A verse in the \textit{Kinsen Roku} runs: \textit{``When the bullish turn appears
and all are bullish --- because the price is high, buying is
forbidden.''} This, too, expresses the same idea as the present
chapter.

% ---------------------------------------------------------------------------
%  Article 17
% ---------------------------------------------------------------------------
\article{17}{Do Not Buy at the Ceiling; Do Not Sell at the Floor}

\begin{sokyubox}
{\small\itshape\color{sokyubar} \jp{第十七章\quad 天井を買はず底を賣らずと申事}}

\medskip
Do not buy at the ceiling (\jt{天井}{tenj\=o}); do not sell at the floor
(\jt{底}{soko}).

This is the first principle. Whether the market is rising or falling,
one does not know where the ceiling or the floor lies. Having no
sense of where the top or the bottom may be, one chases the rally or
presses the decline, and in the end suffers losses. When the market
has risen too far, understand that it will assuredly fall thereafter.
When it is falling, understand that it will assuredly rise. At such
moments, let go of greed and form one's assessment.
\end{sokyubox}

\commentary
When one's forecast has been correct for several months running, and
the market keeps rising or keeps falling, one loses sight of where
the ceiling or the floor might be. One imagines it will go on rising
forever, or falling without limit, and forgets to gauge the position
of the ceiling and the floor. Carried away by momentum, one piles on
buy orders at the ceiling or sell orders at the floor, and the result
is loss. Therefore, when the ceiling draws near, one must never let
go of the determination to sell. When one senses the floor is near,
one must never forget the resolve to buy. The \textit{Yagi Tora no Maki}
admonishes: ``Do not sell at the ceiling; do not buy at the floor.''
The meaning there differs slightly from the present chapter: it
teaches that the ceiling and the floor are never immediately
recognizable, so one should confirm them thoroughly before trading.

% ---------------------------------------------------------------------------
%  Article 18
% ---------------------------------------------------------------------------
\article{18}{Famine in a Bountiful Year; Bounty in a Famine Year}

\begin{sokyubox}
{\small\itshape\color{sokyubar} \jp{第十八章\quad 豊年の凶作凶作の豐年の事}}

\medskip
There is a thing called ``famine in a bountiful year'' and ``bounty in
a famine year.''

When two or three years of good harvests have piled rice high in
every province, a single poor crop does not cause serious concern.
Conversely, when four or five years of half-yields at sixty or
seventy percent have left rice in short supply, even a good harvest
will send prices higher. In general, when good harvests have
continued for two or three years, people forget the feeling of
scarcity; they grow complacent, spending freely without restraint.
At such times, understand that within a year or two, prices may
reverse themselves as if by the flip of a hand.
\end{sokyubox}

\commentary
The first half of this chapter explains ``bounty in a famine year,''
and the second half comments on ``famine in a bountiful year.''

The volume of rice carried forward through reduced consumption is by
no means small. Yet in general, when good harvests continue for two
or three years, people forget the feeling of scarcity. Consumers ---
in weddings and funerals, in celebrations, and in daily life alike ---
all tend toward rice-eating, and this increase in demand is also by
no means negligible. Consequently, in such years, within a year or
two the supply of rice reverses into shortage, and unexpectedly high
prices may appear.

In recent times, when the domestic harvest fails, foreign rice can be
imported to make up the shortfall in supply. Even so, the price of
rice rises in proportion to the domestic crop failure. One should
keep this chapter in mind and maneuver (\jt{駈引}{kakehiki}) accordingly.

Now, when three, four, or even five years of abundant harvests have
continued, and surplus rice from deferred consumption has accumulated,
a single year of poor harvest will not deeply affect supply and
demand, because the surplus compensates for the shortage. Conversely,
after a run of poor harvests, even a single good year will send rice
prices higher. Moreover, when the populace is prosperous, the
standard of living naturally tends toward luxury; ordinary meals shift
from mixed grains toward pure rice. But when harvests have been poor
or the economy depressed, the reverse occurs: even those who formerly
ate only rice switch to barley. The volume represented by this shift
in eating habits is by no means trivial.

% ---------------------------------------------------------------------------
%  Article 19
% ---------------------------------------------------------------------------
\article{19}{What Seems Scarce Will Prove Surplus; What Seems Surplus Will Prove Scarce}

\begin{sokyubox}
{\small\itshape\color{sokyubar} \jp{第十九章\quad 足らぬは餘る餘るは足らぬと申事}}

\medskip
It is said: what seems scarce will prove surplus; what seems surplus
will prove scarce.

When something is plentiful, everyone assumes there is plenty and
lets down their guard, making no preparations. For that very reason,
it ultimately falls short. When something is scarce, everyone is
vigilant and takes pains to secure it. For that very reason, it
ultimately proves more than sufficient.
\end{sokyubox}

\commentary
In a year of poor harvest, everyone is on guard. People develop the
urge to stockpile even rice they have no immediate need for; buy
speculators appear; and farming households hold back their grain,
reluctant to let it go. The market supply consequently falls short.
But when the following spring arrives and the wheat crop looks
promising and the weather normalizes, sellers gradually appear and
physical rice flows into the market. Then, if the seventh and eighth
months bring strong sunshine and prospects for a good harvest firm up,
spot rice floods the market all at once and momentarily piles up.
See also Upper Volume, Article~55, on the sixth-month collapse.

% ---------------------------------------------------------------------------
%  Article 20
% ---------------------------------------------------------------------------
\article{20}{Do Not Envy Another's Trade}

\begin{sokyubox}
{\small\itshape\color{sokyubar} \jp{第二十章\quad 人の商を羨むべからざる事}}

\medskip
Do not look enviously upon another's trade.

The \textit{Yagi Tora no Maki} states: ``Trading born of envy is a thing to
be guarded against.'' This is the meaning of the present chapter.
\end{sokyubox}

% ---------------------------------------------------------------------------
%  Article 21
% ---------------------------------------------------------------------------
\article{21}{One Must Not Trade in Anger}

\begin{sokyubox}
{\small\itshape\color{sokyubar} \jp{第二十一章\quad 腹立売買不可致事}}

\medskip
Selling in anger, buying in anger --- one must never do these things.
Exercise the greatest caution.
\end{sokyubox}

\commentary
To form one's market assessment, one must rely on reason; one must
never let emotion take over. This has been stated repeatedly in
preceding chapters. To trade out of envy for another's success, or
to trade recklessly because a string of losses has driven one to
fury --- all of this means acting not from a reasoned assessment of
the market but from momentary emotion. One must exercise the greatest
caution.

% ---------------------------------------------------------------------------
%  Article 22
% ---------------------------------------------------------------------------
\article{22}{Ceiling and Floor Prices Continue for Three Years, Then Change}

\begin{sokyubox}
{\small\itshape\color{sokyubar} \jp{第二十二章\quad 天井直段底直段三年続き夫より變ずる事}}

\medskip
Ceiling prices and floor prices continue for three years, then change.
Of this there is no doubt.

\begin{sloppypar}
Verify this against the records from Meiwa~5 (\jp{明和五年}, 1768) and it
holds true without exception. First, in the year of a three-year
block (\jt{三年塞}{sannen-fusagari}), when the seventh month falls on
\textit{kinoe} (\jp{甲}), the pattern especially continues for three years. In
other years too, when there is extreme crop failure, it continues for
three years. Observe carefully. Consult the price-rating records to
confirm. If in the first year the market rises 120 \jt{俵}{tawara},
in the second year expect a rise of 100 tawara. When
the market declines from the following spring through summer, the
same ordering applies: the decline is less in the second and third
years than in the first. This also depends on the volume of buy
positions being carried (\jt{買ハラミ}{kai-harami}). Consider well.
\end{sloppypar}
\end{sokyubox}

\commentary
The meaning of this chapter is clear enough from the proviso, but
what remains unclear is the length of one full cycle.

In the \textit{Sanmi Mond\=o} (\jp{三位問答}) by the elder Kuzuoka
(\jp{葛岡翁}), a high-price cycle lasts three years, a mid-price
cycle three years, and a low-price cycle three years, after which it
returns to the high-price cycle. This theory coincides with the view
current in Western commercial society that panics (\jt{恐慌}{ky\=ok\=o})
occur once every ten years.

Testing the present chapter's thesis against Meiji~3, Meiji~13, and
Meiji~23, the pattern appears to hold in every case. S\=oky\=u, however,
regarded the years when the seventh month falls on \textit{kinoe} during a
three-year block facing west as the most reliable instances.

% ---------------------------------------------------------------------------
%  Article 23
% ---------------------------------------------------------------------------
\article{23}{The Analogy of the Three Strategies and the Six Teachings}

\begin{sokyubox}
{\small\itshape\color{sokyubar} \jp{第二十三章\quad 三略六韜タトへの事}}

\medskip
The \textit{Three Strategies} and the \textit{Six Teachings}
(\jt{三略六韜}{Sanryaku Rikut\=o}) are writings on
the innermost secrets of military arts and strategy --- techniques for
fortifying defenses
and forming battle arrays. Yet even with them, one cannot defeat the
enemy and claim victory every time. The rice trade is the same as
military strategy. Though tens of thousands of people trade, none
possess a proper method. This Three-Position Tradition
(\jt{三位の傳}{sanmi no den}) has been devised with even greater freedom
than the \textit{Sanryaku} and \textit{Rikut\=o}. The seventh month turning
to \textit{kinoe}, the three-year convergence --- this is precisely the
Eight Formations (\jt{八陣}{hachijin}). Treasure it and keep it with
the greatest care.
\end{sokyubox}

\commentary
The \textit{Three Strategies} (\jt{三略}{Sanryaku}) and the \textit{Six Teachings}
(\jt{六韜}{Rikut\=o}) are both titles of Chinese military classics. The
\textit{Sanryaku} is divided into Upper, Middle, and Lower Strategies and
treats principally the general's proper conduct. The \textit{Rikut\=o} is
divided into six sections --- Civil, Martial, Tiger, Leopard, Dragon,
and Dog --- and records the maneuvering (\jt{駈引}{kakehiki}) of armies
in the field.

The meaning is this: the \textit{Sanryaku} and \textit{Rikut\=o} expound the
innermost secrets of strategy, yet even with them one cannot win a
hundred battles out of a hundred. Now, the maneuvering of the rice
trade is of the same nature as the unorthodox and orthodox
maneuvers of war. Tens of thousands of people engage in it, yet most
do not know the methods and tactical play of the rice trade. This
book surpasses even the \textit{Sanryaku} and \textit{Rikut\=o} in revealing the
secrets of rice-trade maneuvering; if one follows its teachings, one
will never suffer loss. For example, setting a buy net when the
seventh month turns to \textit{kinoe} during a three-year block is like the
Eight Formations of military science, which let no enemy escape --- a
secret method yielding a hundred gains without a single loss. Keep
this deeply in mind.

% ---------------------------------------------------------------------------
%  Article 24
% ---------------------------------------------------------------------------
\article{24}{When the Seventh Month Turns to Kinoe}

\begin{sokyubox}
{\small\itshape\color{sokyubar} \jp{第二十四章\quad 七月甲に廻る時の事}}

\medskip
When the seventh month turns to \textit{kinoe} (\jp{甲}, the first Heavenly
Stem), both locally and in the Kamigata (\jt{上方}{Kamigata}, the
Osaka--Kyoto region), it is usually a year of natural disaster. At
such times, the market rises sharply from the seventh month; fifty
to sixty tawara in the seventh and eighth months, and by the ninth,
tenth, and November (\jt{霜月}{shimotsuki}) a ceiling of one hundred
to one hundred twenty tawara or more will appear.

On the other hand, when both Kamigata and locally the weather has
been fine from the mid-summer \jt{土用}{doy\=o} onward, and reports of
a good local harvest circulate, and Kamigata too has favorable
conditions, and moreover old rice remains through the fifth month
with roughly half still in the government storehouses --- in a year
when both new and old rice carry over, there will be no
rise through the seventh month. In such a year, the November contract
(\jt{霜月限}{shimotsuki-giri}) will rise gradually and naturally from
the eighth, ninth, and tenth months, and the ceiling price will
appear by the twelfth month or the first month. Do not rush from the
sixth month onward. Wait and observe. Consider well: the volume of
old rice, the quality of the harvest, and the state of market
sentiment (\jt{人氣}{ninki}) all differ. For this reason no two years
are alike. Yet whether early or late, in such a year the market will
invariably rise.
\end{sokyubox}

\commentary
The maneuvering for a year when the seventh month turns to \textit{kinoe}
is explained in detail in Upper Volume, Article~51. Here it is
stated more briefly.

Now, when the seventh month turns to \textit{kinoe} and natural disaster
strikes, the market rises steadily for five or six months from around
the seventh month. However, if that year's weather is favorable and
stored rice is plentiful, the rise is delayed, beginning around the
eighth or ninth month. Still, because the volume of old rice, the
quality of the harvest, and the state of sentiment cause the timing
to vary, one must weigh these factors and maneuver accordingly.

% ---------------------------------------------------------------------------
%  Article 25
% ---------------------------------------------------------------------------
\article{25}{On the Kamigata Price Scale}

\begin{sokyubox}
{\small\itshape\color{sokyubar} \jp{第二十五章\quad 上方相塲位附の事}}

\medskip
\begin{itemize}[nosep]
\item Chikuzen rice (\jp{筑前米}) trades about 4 \textit{monme} below Higo rice (\jp{肥後米}).
\item Kaga rice (\jp{加賀米}) trades about 5 \textit{monme} below Chikuzen rice.
\item Ch\=ugoku rice (\jp{中國米}) trades about 9 \textit{monme} below the same [Chikuzen] rice.
\item Sh\=onai rice (\jp{莊内米}) trades about 7 \textit{monme} below the same [Chikuzen] rice.
\end{itemize}

\medskip
The above proportions have held year after year.
\end{sokyubox}

\commentary
The four items above are the price-rating table (\jt{格附表}{kakutsuke-hy\=o})
as it stood in Osaka at that time.

% ---------------------------------------------------------------------------
%  Article 26
% ---------------------------------------------------------------------------
\article{26}{On Trading by the Contract Month}

\begin{sokyubox}
{\small\itshape\color{sokyubar} \jp{第二十六章\quad 月切商仕様の事}}

\medskip
In a one-month contract trade (\jt{月切}{tsukigiri}), one should gauge
the moment and take even a small profit.
\end{sokyubox}

\commentary
When a position (\jt{商内}{akinai}) one has established reaches the
current contract month's expiry, even if one still sees some
prospect ahead, one should generally cut the position and secure
whatever small profit remains.

% ---------------------------------------------------------------------------
%  Article 27
% ---------------------------------------------------------------------------
\article{27}{On the Doyo Entry and Seasonal Turns: The Yin-Yang Counting Method}

\begin{sokyubox}
{\small\itshape\color{sokyubar} \jp{第二十七章\quad 土用入并に節替り陰陽繰方の事}}

\medskip
The method for knowing whether one enters yin (\jp{陰}) or yang
(\jp{陽}) at the \jt{土用入}{doy\=o-iri} and at the
seasonal turn (\jt{節替り}{setsu-gawari}):

\medskip
\begin{small}
\begin{tabular}{@{}p{3cm}p{6.5cm}@{}}
\toprule
Odd months (1, 3, 5, 7, 9, 11) & \jp{子寅辰午申戌} (\textit{Ne, Tora, Tatsu, Uma, Saru, Inu}) \\[0.3em]
Even months (2, 4, 6, 8, 10, 12) & \jp{丑卯巳未酉亥} (\textit{Ushi, U, Mi, Hitsuji, Tori, I}) \\
\bottomrule
\end{tabular}
\end{small}

\medskip
Count from the first day onward to determine the day. When the
entry falls on a Day of the Rat (\jp{子}), regard it as yang; when
it falls on a Day of the Ox (\jp{丑}), regard it as yin. The
remaining days follow this pattern accordingly.
\end{sokyubox}

% ---------------------------------------------------------------------------
%  Article 28
% ---------------------------------------------------------------------------
\article{28}{On Kinoe, Kinoto, and Unlucky Days}

\begin{sokyubox}
{\small\itshape\color{sokyubar} \jp{第二十八章\quad 甲乙並に不成就日の事}}

\medskip
The ten Heavenly Stems (\jt{十干}{jikkan}): \jp{甲} (\textit{kinoe}),
\jp{乙} (\textit{kinoto}), \jp{丙} (\textit{hinoe}),
\jp{丁} (\textit{hinoto}), \jp{戊} (\textit{tsuchinoe}),
\jp{己} (\textit{tsuchinoto}), \jp{庚} (\textit{kanoe}),
\jp{辛} (\textit{kanoto}), \jp{壬} (\textit{mizunoe}),
\jp{癸} (\textit{mizunoto}).

\medskip
Unlucky days (\jt{不成就日}{f\=uj\=ojubi}):

\begin{itemize}[nosep]
\item \textbf{1st and 7th months}: 3rd, 11th, 19th, 27th
\item \textbf{2nd and 8th months}: 2nd, 10th, 18th, 26th
\item \textbf{3rd and 9th months}: 1st, 9th, 17th, 25th
\item \textbf{4th and 10th months}: 4th, 12th, 20th, 28th
\item \textbf{5th and 11th months}: 5th, 13th, 21st, 29th
\item \textbf{6th and 12th months}: 6th, 14th, 22nd, 30th
\end{itemize}
\end{sokyubox}

% ---------------------------------------------------------------------------
%  Article 29
% ---------------------------------------------------------------------------
\article{29}{On the Sixth-Month Doyo Entry and the Day of the Ox}

\begin{sokyubox}
{\small\itshape\color{sokyubar} \jp{第二十九章\quad 六月土用入丑の日の事}}

\medskip
If within the first six days after the entry into the sixth-month
\jt{土用}{doy\=o}, a Day of the Ox (\jt{丑の日}{ushi no hi}) falls,
autumn winds arrive early and the harvest suffers. However, if it
falls on the second or third day, the damage is particularly severe;
if it falls on the fourth, fifth, or sixth day, the effect is not
so great.

Furthermore, if it rains on the fifth day after the \textit{doy\=o} entry,
then even if the crop inspection looks favorable, the actual yield
will fall short. When the Day of the Ox falls on the third day, it
is called ``Doy\=o Sabur\=o'' (\jp{土用三郎}). In such a year, moreover,
price fluctuations are pronounced.

If a Day of the Tiger (\jt{寅の日}{tora no hi}) falls within the first
three days, this too is unfavorable. This is because when the Day
of the Ox falls on the second day, the Day of the Tiger necessarily
falls on the third. Counting five days from each \textit{doy\=o} entry, the
autumn period amounts to approximately thirty days; hence these are
dates of critical importance.
\end{sokyubox}

% ---------------------------------------------------------------------------
%  Article 30
% ---------------------------------------------------------------------------
\article{30}{On Rain on the Seventh Day of the Seventh Month}

\begin{sokyubox}
{\small\itshape\color{sokyubar} \jp{第三十章\quad 七月七日雨の事}}

\medskip
Rain on the seventh day of the seventh month (\jt{七夕}{Tanabata}):
heavy rain does not signify much, but light rain is decidedly bad
for the harvest. This holds without exception.
\end{sokyubox}

% ---------------------------------------------------------------------------
%  Article 31
% ---------------------------------------------------------------------------
\article{31}{On the Autumn Equinox}

\begin{sokyubox}
{\small\itshape\color{sokyubar} \jp{第三十一章\quad 秋の彼岸の事}}

\medskip
When the autumn equinox (\jt{秋の彼岸}{aki no higan}) extends even a
single day into the ninth month, the harvest will without question
be poor. This rule has never failed; the reason must be that the
season runs late. Even if the equinox does not reach into the ninth
month, if it falls after the twentieth of the eighth month, the
outcome is unfavorable.

Furthermore: if a Day of the Rooster (\jt{酉の日}{tori no hi}) falls
within the seven days of the equinox period, a great wind will
blow without fail.
\end{sokyubox}

% ---------------------------------------------------------------------------
%  Article 32
% ---------------------------------------------------------------------------
\article{32}{On Solar Eclipses}

\begin{sokyubox}
{\small\itshape\color{sokyubar} \jp{第三十二章\quad 日蝕の事}}

\medskip
A solar eclipse (\jt{日蝕}{nisshoku}) occurring between the hour of
\textit{yotsu-mae} (\jp{四つ前}, approximately 9~a.m.) and the hour past
\textit{yatsu} (\jp{八つ過}, approximately 3~p.m.) is bad for the harvest;
eclipses outside these hours are said not to matter. Nevertheless,
any year in which either a solar or lunar eclipse occurs has without
exception, whether great or small, been a year of poor harvest.
\end{sokyubox}

% ---------------------------------------------------------------------------
%  Article 33
% ---------------------------------------------------------------------------
\article{33}{On Unlucky Days}

\begin{sokyubox}
{\small\itshape\color{sokyubar} \jp{第三十三章\quad 厄日の事}}

\medskip
The unlucky days (\jt{厄日}{yakubi}): it is said that if the
eighty-eighth night (\jt{八十八夜}{hachij\=uhachiya}) passes without
incident, then the two-hundred-tenth day (\jt{二百十日}{nihyaku t\=oka})
will also pass without incident.

However, even if the two-hundred-tenth day passes safely, should a
Day of the Rooster fall on the third day thereafter, a great storm
will come on that day. This rule holds without exception. In Kansei~9
(\jp{寛政九年}, 1797), the Day of the Rooster fell on the third day,
and indeed a great storm struck.
\end{sokyubox}

% ---------------------------------------------------------------------------
%  Article 34
% ---------------------------------------------------------------------------
\article{34}{On the Three Trees}

\begin{sokyubox}
{\small\itshape\color{sokyubar} \jp{第三十四章\quad 三木の事}}

\medskip
\begin{itemize}[nosep]
\item When \jp{甲乙} (the Wood stems) continue for three months running, great winds will blow.
\item When \jp{丙丁} (the Fire stems) continue for three months, there will be drought.
\item When \jp{戊己} (the Earth stems) continue for three months, know that the price of grain will be high.
\item When \jp{庚辛} (the Metal stems) continue for three months, there is reason to expect a grain shortage.
\item When \jp{壬癸} (the Water stems) continue for three months, grain prices rise high and people perish.
\end{itemize}

\medskip
Consider the above in relation to the season --- spring, summer,
autumn, or winter --- and weigh the implications for great winds,
grain prices, drought, and so forth.

In Tenmei~6 (\jp{天明六年}, 1786), the first of the ninth month fell on
\jp{庚丑} (\textit{kanoe-ushi}), the first of the tenth month on
\jp{庚未} (\textit{kanoe-hitsuji}), and the first of the intercalary tenth
month on \jp{辛未} (\textit{kanoto-hitsuji}) --- the Metal stems continuing
for three months in succession. That year, the price stood at
26~tawara 8~\textit{bu}, reaching a ceiling of 16~tawara 4~\textit{bu} on the
sixteenth of the tenth month. That year, the autumn equinox fell on
the twenty-eighth of the eighth month.
\end{sokyubox}

\commentary
``When \jp{甲乙} continue for three months'' means: when the first
day of the month (\jt{朔日}{tsuitachi}) falls on a \textit{kinoe} or
\textit{kinoto} stem for three consecutive months. The same principle
applies to \jp{丙丁} and all the rest.

% ---------------------------------------------------------------------------
%  Article 35
% ---------------------------------------------------------------------------
\article{35}{On Fire Days and Wood Days Occurring Three or Four in a Row}

\begin{sokyubox}
{\small\itshape\color{sokyubar} \jp{第三十五章\quad 朔日火木三つ四つ續く事}}

\medskip
\begin{itemize}[nosep]
\item When the first of the month falls on a Fire-element day (\jp{火}) four times in succession, that year will see drought.
\item When the first of the month falls on a Wood-element day (\jp{木}) three times in succession, that year will see frequent great storms.
\item In any month, when Water-element days (\jp{水}) occur five times in succession, it is said that ``five waters destroy the seed,'' meaning frequent flooding will occur.
\end{itemize}

\medskip
However, this rule is not rigid --- one should consider it carefully.
\end{sokyubox}

% ---------------------------------------------------------------------------
%  Article 36
% ---------------------------------------------------------------------------
\article{36}{On Days of the Ox and Days of the Dragon}

\begin{sokyubox}
{\small\itshape\color{sokyubar} \jp{第三十六章\quad 丑の日辰の日の事}}

\medskip
Days of the Ox (\jt{丑の日}{ushi no hi}) and Days of the Dragon
(\jt{辰の日}{tatsu no hi}) are days when rice is cheap. However, the
Day of the Ox is also a day on which floor prices and ceiling prices
tend to emerge. At ceiling and floor turning points, one should pay
close attention.
\end{sokyubox}

% ---------------------------------------------------------------------------
%  Article 37
% ---------------------------------------------------------------------------
\article{37}{On the Saying ``One Mark, Then Two''}

\begin{sokyubox}
{\small\itshape\color{sokyubar} \jp{第三十七章\quad 壹つ目二つと申事}}

\medskip
The saying ``one mark, then two, with no rise or fall'' refers to the
``marks'' (\jp{目}, \textit{me}) in the price-ranking ledger
(\jt{位付の書}{kuraizuke no sho}). One should understand this accordingly.
\end{sokyubox}

\commentary
Regarding the ``one mark'' of this chapter, one should refer to Upper
Volume, Article~27, where it states: ``In the first and second months,
if the first month rises then the second month falls; if the first
month falls then the second month rises.''

% ---------------------------------------------------------------------------
%  Article 38
% ---------------------------------------------------------------------------
\article{38}{On the Period from the First through the Fourth or Fifth of the Month}

\begin{sokyubox}
{\small\itshape\color{sokyubar} \jp{第三十八章\quad 朔日より四五日頃迄の事}}

\medskip
When rice dips once or twice between the first and the fourth or
fifth of the month, this constitutes a buying opportunity within a
range-bound market (\jt{通相場}{kayoi s\=oba}). This refers to the pattern
described earlier: during a falling market after a ceiling, when the
beginning of the current month is higher than the end of the previous
month, prices will rise that month. One should take note of this.
\end{sokyubox}

\commentary
As has already been explained, a ``falling market after a ceiling''
is one in which, for a period of two or three months or even five
or six months, prices are high at the start of each month and low at
month-end. However, if this pattern is broken --- if, carrying over
from the end of the previous month, the beginning of the current
month posts a high price --- the market has shifted to an upward
range-bound phase (\textit{kayoi-ue}), and prices will rise through the
month. One should bear this in mind.

% ---------------------------------------------------------------------------
%  Article 39
% ---------------------------------------------------------------------------
\article{39}{Guidance for the Rice Trader}

\begin{sokyubox}
{\small\itshape\color{sokyubar} \jp{第三十九章\quad 米商心得方の事}}

\medskip
The principle of ``at two, close out; at three, fully ripe; at four,
it turns'' (\jp{二つ仕舞、三つ十分、四つ轉じ}) has been set forth
previously.

The rice trade proceeds from two directions only --- rise and fall ---
yet the paths of the human heart are as numerous and tangled as
threads. Those who devote themselves to this trade should constantly
and vigilantly examine the principles set forth in the preceding
chapters. When one has thoroughly discerned the whole, one will
naturally come to know the strength and weakness of the rice market.
When one studies this book carefully and trades accordingly, then no
matter how the market moves --- be it from natural disasters,
unseasonal weather, bountiful or poor harvests, abundance or
scarcity of rice, volume of shipping, or the force of market
sentiment --- one will never fail to profit. This must absolutely not
be shown to others; guard it with the utmost care.
\end{sokyubox}

\commentary
The phrase ``at two, close out; at three, fully ripe; at four, it
turns'' is the supreme principle of futures trading (\jt{定期賣買}{teiki baibai}).
Whether in contract rice, spot rice, or equities, the tactical
maneuvering (\jt{駈引}{kakehiki}) of buying and selling is contained
entirely in this single saying --- it may fairly be called a golden
maxim. Accordingly, its application may be interpreted from various
angles.

Some interpret it in light of the passage elsewhere in this book:
``After the ceiling month, prices will assuredly fall for two or three
months; by the fourth month, depending on the commodity, they may
rise --- one should not sell.'' On this reading, one follows the
market for two or three months, and by the fourth month, the trend
reverses.

Others read it as: take profit at seven or eight parts, close out
the position when it is fully ripe, then turn one's mind --- if one
had been buying, switch to selling; if one had been selling, switch
to buying.

Still others say it describes the rhythm of the market: take profit
on a position after one or two swings, wait for the market to reach
its full extreme, identify the highest or lowest point, and then
play the reversal.

In the end, from whichever angle one interprets it, this is the
innermost secret of market maneuvering --- there is no other
``secret of the market'' beyond this. In particular, the appended
remark --- \textit{``The rice trade proceeds from two directions only, rise
and fall, yet the paths of the heart are as numerous and tangled as
threads''} --- is a golden maxim that spans East and West, past and
present; there is no finer truth than this.

The meaning is this: the rice market is a thing whose causes of rise
and fall are exceedingly complex, and these complex causes act upon
the human heart, manifesting as the market. Because the market is
formed by the hearts of ten thousand people gathered together, the
paths of those hearts are as tangled as threads. Yet in the end, just
as a thread has two ends, there are only two directions --- up and
down --- and nothing else. Put another way, observing the rise and
fall of the market is like trying to untangle knotted thread. When
one harbors feelings of attachment, greed, or suspicion, or when
one's mind grows impatient and one tries to profit in haste, one
cannot arrive at a sound judgment --- yet forces oneself to form one
anyway. In such a state, it is exceedingly difficult to find the two
ends --- the ceiling and the floor. In most cases, what is not the
thread's end appears to be the end, so one opens a position only to
encounter a move in the completely opposite direction and suffer
great loss.

How, then, should one identify those two ends --- the ceiling and
the floor? The method is to study the explanations set forth in
the preceding chapters of this book, reading them over and over
until they are thoroughly absorbed. In the course of reading, one
will naturally come to perceive the strength and weakness of the
market. When one has fully mastered this book and trades on that
basis, then no matter what great swings one encounters, there will
be no confusion, and one will be able to maneuver freely and at
will --- and thereby amass great wealth. This is, without question,
a secret book that must never be shown to others.

% ---------------------------------------------------------------------------
%  Article 40
% ---------------------------------------------------------------------------
\article{40}{Guidance on Weather Conditions}

\begin{sokyubox}
{\small\itshape\color{sokyubar} \jp{第四十章\quad 一天心得方の事}}

\medskip
There is an old saying: ``When the sky is completely clear, a flood
will come within three days.'' This is indeed true. When fine weather
continues for several days, it will without fail be followed by
prolonged rain. Moreover, in the sixth and seventh months, when the
sky is perfectly clear without a wisp of cloud and the heat is
extreme, heavy rain will invariably come within two or three days.
In some years, however, a long dry spell sets in and twenty or
thirty days or more may pass without rain; one should judge according
to the circumstances.
\end{sokyubox}

% ---------------------------------------------------------------------------
%  Article 41
% ---------------------------------------------------------------------------
\article{41}{On Buying and Selling at Ceiling and Floor in the November Contract}

\begin{sokyubox}
{\small\itshape\color{sokyubar} \jp{第四十一章\quad 霜月限天井底賣買の事}}

\medskip
In trading the November contract (\jt{霜月限}{shimotsuki-giri}), over
the course of several months there are shifts in timing, and both
the buyer's conviction and the seller's judgment may prove correct
at different points. For example, even if a buyer enjoys a run of
profitable fortune, should he continue piling on purchases at the
ceiling, his fortune will turn adverse. Likewise, even if a seller
holds losing positions, should he press his selling at the very
point where prices reach the ceiling, his fortune too will turn
adverse. When a rise is pushed to its limit, a fall follows; when
a fall is pushed to its limit, a rise follows. This is the
principle of yin and yang (\jt{陰陽}{in'y\=o}).
\end{sokyubox}

\commentary
The meaning is clear without awaiting commentary. The essential point
is that one must trade the ceiling and the floor.

% ---------------------------------------------------------------------------
%  Article 42
% ---------------------------------------------------------------------------
\article{42}{On Fluctuations in a Range-Bound Market and Cutting One's Losses}

\begin{sokyubox}
{\small\itshape\color{sokyubar} \jp{第四十二章\quad 通の高下見切の事}}

\medskip
In a range-bound market, target a swing of ten \jt{俵}{tawara}
and cut quickly --- this is the first principle.
\end{sokyubox}

\commentary
Before opening a trade, it is essential to determine whether the
market is range-bound (\jt{通相塲}{kayoi s\=oba}) or trending
(\jt{伸相塲}{nobi s\=oba}) and to trade accordingly. Range-bound
fluctuations involve narrow price swings, so one must not employ the
leisurely maneuvering suited to a trending market. One should target
a swing of ten tawara, take profit (\jt{利喰}{rigui}) promptly, and
close the position.

Elsewhere in this book it states: ``In a range-bound market, where
swings are only one or two tawara, one should quickly cut at a ten-
or twenty-tawara swing and secure the profit'' --- the meaning is
identical to this chapter.
